\section* {Лекция 8. Продолжение деревьев решений}

\subsection*{Параметры дерева}
Прежде чем продолжать наши рассуждения, давайте сразу обсудим понятие параметров дерева.

Ранее мы уже сталкивались с \textit{гиперпараметрами} "--- например \texttt{learning rate} в алгоритме градиентного спуска, или коэффициентом регуляризации. Здесь аналогичная история лишь с разницей в том, что у деревьев параметров значительно больше:
\begin{itemize}
  \item \textbf{max\_depth} — максимальная глубина дерева; ограничивает количество уровней ветвления и помогает предотвратить переобучение модели.
  
  \item \textbf{min\_samples\_leaf} — минимальное количество объектов выборки, которые должны находиться в листовой вершине; регулирует сложность дерева и его обобщающую способность.
  
  \item \textbf{criterion} — критерий качества для выбора оптимального разбиения на каждом узле (о которых позже); определяет, насколько хорошо ветвление разделяет классы.
  
  \item \textbf{max\_features} — максимальное количество признаков, рассматриваемых при разбиении узла;

  \item \dots
\end{itemize}

